\documentclass[10pt,a4paper]{amsart}
\usepackage[margin=19mm]{geometry}
\usepackage{amsmath,amssymb,mathtools}
\usepackage[scr=rsfs,cal=euler]{mathalfa}
\usepackage{xfrac}
\usepackage[unicode,hidelinks,bookmarks=false]{hyperref}
\newcommand{\ie}{i.e.\ }
\newcommand{\cf}{cf.\ }
\newcommand{\resp}{respectively\ }
\newcommand{\Subsection}{\subsection}
\providecommand{\nobreakdash}{\nobreak\hbox{-}}
\setlength{\emergencystretch}{2em}
\multlinegap0pt
\usepackage{bm}
\usepackage{bbm}
\usepackage{dsfont}
\usepackage{xspace}
\usepackage{cancel}
\usepackage{mathtools}
\usepackage{amsthm}
\newtheorem{theorem}{Theorem}
\newtheorem{lemma}[theorem]{\bf Lemma}
\newcommand{\pr }{\mathrm{Pr} }
\title[Coprime automorphisms of profinite groups and the commuting ]{Coprime automorphisms of profinite groups and the commuting probability of Sylow subgroups}
\author{Eloisa Detomi, Robert M. Guralnick, Marta Morigi, Pavel Shumyatsky}
\date{Source: arXiv:2609.29925v1 (CC BY 4.0)}
\begin{document}
\maketitle

\noindent\emph{Source and licence.} Coprime automorphisms of profinite groups and the commuting probability of Sylow subgroups. Version arXiv:2609.29925v1 (CC BY 4.0), distributed under the Creative Commons Attribution 4.0 International licence, as recorded in the arXiv licence metadata for this exact version and shown on the abstract page https://arxiv.org/abs/2609.29925v1. Full rights notice: licenses/arxiv-2609.29925-CC-BY-4.0.txt.\par\smallskip

\noindent\textit{Editorial scope.} Counted unit: the Lemma whose source label is \texttt{prod} in the article (source0.tex lines 222--225, source label \texttt{prod}), delivered together with the article's own complete original proof of it (source0.tex lines 226--257, one \texttt{proof} environment of 32 lines). No mathematics is written, repaired or re-derived: statement and proof are the article's own text line for line. Required context retained uncounted: quot (lines 208--211); limit (lines 217--220). Every definition, display and label of those lines is retained unchanged. Omitted from this package and not redistributed: 1--207, 212--216, 221--221, 258--852. Nothing inside a retained excerpt is omitted or altered: every retained body is delivered exactly as the article's lines are written, so no omission receipt of any kind accompanies this package. Citations: the retained excerpts carry no citation command at all, so no bibliography entry of the article is transcribed and no reference is redistributed. Reference adaptation: none was needed and none was made; every reference inside the delivered unit resolves inside it, so no unresolved reference or citation is produced. Printed slips kept literal: no printed word, symbol, hypothesis or number of the counted statement or of its proof was altered. Layout and class: the article's own class is \texttt{amsart} with packages including amsmath, amssymb, enumerate, color; the wrapper is the lane's pinned \texttt{amsart} editorial wrapper at 10pt on A4, which restates the theorem-like environments the retained text needs and adds the article's own macro definitions that the retained text needs (\texttt{\textbackslash pr}), with extra wrapper packages (bm, bbm, dsfont, xspace, cancel, mathtools). The delivered numbering is the unit's own. Assets: no retained span contains a figure environment, a graphics-inclusion command, a PDF-page inclusion, a table or a TikZ picture; no image file is copied into this package, the package declares no asset dependency and no image file is redistributed with it. Licence: version arXiv:2609.29925v1 of this article is distributed under the Creative Commons Attribution 4.0 International licence, as recorded in the arXiv licence metadata for this exact version and shown on the abstract page https://arxiv.org/abs/2609.29925v1. A full rights notice is delivered at licenses/arxiv-2609.29925-CC-BY-4.0.txt. Characters: every retained excerpt is pure ASCII, so the delivered engine, XeLaTeX with its default Latin Modern text font, reports no missing character.
\begin{lemma}\label{quot} 
 Let $G$ be a compact group and let $N$ be a normal subgroup of  $G$.  For any subgroups $H,K\leq G$ we have 
\[\pr(H,K)\leq \pr(HN/N,KN/N)\pr(N\cap H,N\cap K).\]
\end{lemma}

\begin{lemma}\label{limit} 
Let $G$ be a profinite group and $H, K \le G$.  Then 
\[ \pr(H,K) = \inf_{N \lhd_o G} \pr \left(\frac{HN}{N},\frac{KN}{N}\right).\]
\end{lemma}\

\begin{lemma}\label{prod} 
Let $G$ be a profinite group and $H, K \le G$. Assume that $H=H_1H_2$ is the product of two subgroups. Then
\[\pr(H,K)\ge\pr(H_1,K)\pr(H_2,K)\]
\end{lemma}

\begin{proof} First assume that $G$ is finite. We will repeatedly use tha fact that, if $A$ and $B$ are subgroup of $G$ then every element of the product $AB$ % can be writte in exactly $|H_1\cap H_2|$ ways as a product $h=h_1h_2$, with $h_1\in H_1,h_2\in H_2$. In particular, $|H|= (|H_1|\,|H_2|)/|H_1\cap H_2|$. 
can be written in exactly $|A\cap B|$ ways as a product $ab$, with $a\in A$, $b\in B$. In particular, $|AB|= (|A|\,|B|)/|A\cap B|$. 
We have that
\begin{eqnarray*}
\pr(H,K)
&=&
\frac{1}{|H||K|}
\sum_{h\in H} |C_K(h)|
 = %&=&
\frac{|H_1\cap H_2|}{|H_1||H_2||K|}\sum_{h_1h_2\in H} |C_K(h_1h_2)|\\
&=&
\frac{1}{|H_1||H_2||K|}\sum_{h_1\in H_1}\sum_{h_2\in H_2}
|C_K(h_1h_2)|.
\end{eqnarray*}
%\geq \frac{1}{|H_1||H_2||K|}\sum_{h_1\in H_1}\sum_{h_2\in H_2}|C_K(h_1)\cap C_K(h_2)|=\frac{1}{|H_1||H_2||K|}\sum_{h_1\in H_1}\sum_{h_2\in H_2}
%\frac{|C_K(h_1)|\,|C_K(h_2)|}{|C_K(h_1)C_K(h_2)|}\ge \frac{1}{|H_1||H_2||K|}\sum_{h_1\in H_1}\sum_{h_2\in H_2}\frac{|C_K(h_1)|\,|C_K(h_2)|}{|K|}
 Note that 
 \begin{eqnarray*}
|C_K(h_1h_2)| &\ge& |C_K(h_1)\cap C_K(h_2)|=\frac{|C_K(h_1)|\,|C_K(h_2)|}{|C_K(h_1)C_K(h_2)|} \\
&\ge& \frac{|C_K(h_1)|\,|C_K(h_2)|}{|K|}.
\end{eqnarray*}

Therefore 
 \begin{eqnarray*}
\pr(H,K) &\ge &\sum_{h_1\in H_1}\sum_{h_2\in H_2}\frac{|C_K(h_1)|\,|C_K(h_2)|}{|H_1||H_2||K|^2}\\
 &=& \left(\sum_{h_1\in H_1}\frac{|C_K(h_1)|}{|H_1||K|}\right)\,\left(\sum_{h_2\in H_2}\frac{|C_K(h_2)|}{|H_2||K|}\right) \\
 &=& \pr(H_1,K)\pr(H_2,K).
\end{eqnarray*}

If $G$ is a profinite group, the result follows from Lemma \ref{limit}, taking into account the fact that, by Lemma \ref{quot}, the commuting 
probability of two subgroups does not decrease when considering homomorphic images of $G$.
\end{proof}

\end{document}
